Este trabalho tem como objetivo explorar o Ray Marching, um método menos difundido quando comparado ao Ray Tracing, mas que ainda apresenta muitos casos de uso relevantes. São discutidas suas principais características a partir da implementação de um sistema de renderização de nuvens, destacando as ferramentas oferecidas por essa técnica. ara fundamentar o método, são analisadas as Signed Distance Functions (SDFs), que permitem a modelagem de formas orgânicas e complexas como a das nuvens, que seriam difíceis de reproduzir com as técnicas tradicionais de renderização que envolvem polígonos. Além disso, também é apresentada uma implementação em GLSL de Ray Marching e realizada uma comparação com o Ray Tracing, evidenciando diferenças em termos de desempenho, fidelidade visual e aplicações práticas.